% ----------------------------------------------------------
% Considerações Finais
% ----------------------------------------------------------

\chapter[Conclusão]{Conclusão}

Este trabalho desenvolveu o MetaRefino, aplicativo para Windows que assiste o pesquisador no refino quimiométrico de matrizes de metabolômica não direcionada obtidas por LC-MS a partir das matrizes exportadas pelo Progenesis QI. O aplicativo reúne, num mesmo fluxo, a validação da entrada, o diagnóstico da deriva instrumental, a correção do sinal, o pré-processamento, as análises estatísticas e a exportação no formato do MetaboAnalyst ou CSV, e registra a configuração usada e a origem de cada escolha. O problema que motivou o trabalho era a fragmentação dessas etapas em ferramentas separadas, com escolhas sem critério registrado; o resultado é uma ferramenta que integra as etapas e torna auditáveis as decisões tomadas sobre a matriz, sem retirar do pesquisador a decisão final.

\section{Conclusão dos objetivos}

\textbf{Sistematizar as etapas de pré-processamento.} As etapas do refino foram organizadas numa ordem fixa e declarada: a correção do sinal vem primeiro, como etapa própria, e é seguida pela filtragem por RSD e por valores ausentes, pela imputação, pela normalização, pela transformação e pelo escalonamento. Cada etapa e seus parâmetros ficam registrados com o resultado (Seção~\ref{subsec:pre-processamento-dados}). O objetivo foi atendido.

\textbf{Estruturar um fluxo de trabalho assistido.} O aplicativo recebe o par de arquivos exportado pelo Progenesis QI e conduz o tratamento até a matriz de saída, sem que o pesquisador precise alternar entre ferramentas. Na avaliação com usuários do laboratório parceiro, os relatórios descrevem a importação, a correção, o pré-processamento e a geração das matrizes de saída para o MetaboAnalyst; um deles registra que a correção e a gravação dos arquivos foram feitas sem auxílio (Seção~\ref{sec:5-4}). O objetivo foi atendido, com a ressalva de que essa avaliação é um relato escrito e não observado.

\textbf{Incorporar algoritmos como sistema de apoio à decisão.} O objetivo foi atendido em parte. O diagnóstico classifica a forma da deriva por uma cascata de regras e propõe as fronteiras de lote antes da correção, com justificativa; sobre matrizes sintéticas, recuperou a forma do defeito em 69,5\% dos casos com intensidade moderada, abaixo do alvo de mais de 85\%, e em até 88,5\% com deriva forte. Para a escolha do corretor, um modelo de aprendizado de máquina não superou uma regra fixa que escolhe sempre o SERRF (67,5\% contra 75,0\%, deixando um modo de fora) e deu lugar, na interface, a uma medição do RSD dos QCs fora da amostra numa amostra de 200 \textit{features} da própria matriz, que concordou com a medição completa em 93,1\% dos sorteios simulados. As configurações das etapas não são escolhidas automaticamente: seguem os padrões declarados da ferramenta ou a escolha do pesquisador.

\textbf{Implementar análises estatísticas com critérios de validação.} O aplicativo oferece comparação univariada com correção de Benjamini-Hochberg, PCA, PLS-DA, classificadores por floresta aleatória e por máquina de vetores de suporte e PERMANOVA. O PLS-DA não é apresentado sem o Q$^{2}$ e o \textit{p}-valor de permutação, e os classificadores são acompanhados da acurácia balanceada, do nível do acaso e do \textit{p}-valor de permutação. Com rótulos aleatórios, o \textit{p}-valor de permutação do PLS-DA ficou acima de 0,05 em 95\% das execuções, enquanto o Q$^{2}$ saiu positivo em 14\% delas: o indicador calibrado contra falso positivo é o \textit{p}-valor, e um Q$^{2}$ ligeiramente positivo não basta para afirmar separação entre grupos (Seção~\ref{sec:5-3-3}). O objetivo foi atendido.

\textbf{Avaliar o desempenho em dados reais e sintéticos.} A ferramenta foi avaliada com os dados de um único estudo do laboratório parceiro, em quatro modos cromatográficos, com matrizes sintéticas derivadas deles e com a devolutiva dos usuários. O caminho sem correção reproduz o pipeline em R usado anteriormente pelo laboratório, e a correção SERRF concorda estatisticamente com ele em três dos quatro modos (Seção~\ref{sec:5-3-1}). O resultado central do trabalho veio da comparação entre réguas: medida dentro da amostra, a correção LOESS fica à frente da ausência de correção nos quatro modos; medida fora da amostra, fica atrás dela em três (Seção~\ref{sec:5-3-2}). A régua mais usada para avaliar uma correção favorece, portanto, o método que se ajusta aos próprios QCs, e a ferramenta passou a medir fora da amostra. O objetivo foi atendido, no recorte de um estudo.

\section{Contribuições}

A principal contribuição é a ferramenta: um fluxo único, com as escolhas registradas e justificadas, entregue como executável que não exige instalação do ambiente de desenvolvimento. A segunda é de método: a avaliação da correção por QC fora da amostra, que muda o veredito sobre o melhor corretor em relação à medida usual.

\section{Limitações}

Toda a evidência sobre dados reais vem de um único estudo, os limiares do diagnóstico foram calibrados sobre essas corridas, e as matrizes sintéticas derivam das mesmas. O que funciona nesse espectro pode não funcionar em outros estudos, laboratórios ou equipamentos. A escolha do corretor feita apenas pelos QCs, inclusive a do fluxo anterior do laboratório, acertou o melhor corretor para as amostras em pouco mais de um terço das matrizes sintéticas com deriva gradual, como discutido na seção sobre inteligência artificial e decisões automatizadas, e por isso a decisão final continua com o pesquisador. Funcionalidades de maior apoio aos pesquisadores começaram a ser desenvolvidas, mas precisaram ser classificadas como trabalhos futuros devido ao tempo limitado de uma pesquisa de conclusão de curso. Por fim, a verificação automatizada tem limites próprios, descritos na Seção~\ref{sec:5-5}.

\section{Trabalhos futuros}

Esta monografia descreve o aplicativo nas versões que cada capítulo nomeia. O desenvolvimento continuou depois delas, em grande parte a partir da devolutiva dos técnicos do laboratório, e as funções acrescentadas não estão descritas nos capítulos de desenvolvimento e de testes nem foram avaliadas pelos usuários. Já implementadas no código, numa versão ainda não publicada, e por avaliar:

\begin{itemize}
\item a reorganização da interface em abas por etapa (preparação do estudo, correção, normalização, análise estatística e execução e exportação), com tema claro e escuro;
\item a exportação das figuras em PNG, TIFF ou PDF, com resolução e tamanho escolhidos, e a escolha da cor e do símbolo de cada grupo;
\item a avaliação da normalidade por variável e por grupo depois da correção e da normalização;
\item o intervalo de confiança do \textit{fold change} na análise univariada e nas figuras;
\item as regiões de predição nas PCAs, a PCA colorida por grupo e o gráfico do PLS-DA com os escores de validação cruzada;
\item a curva do LOESS ao longo das injeções, com a sensibilidade à retirada de cada QC, e um intervalo para o RSD de cada corretor;
\item o resumo metodológico e as referências dos métodos usados no processamento exportado;
\item explicações mais claras do critério da sugestão de corretor e um acompanhamento mais detalhado dos trabalhos em execução.
\end{itemize}

Ainda não implementados, e decorrentes dos limites acima:

\begin{itemize}
\item validar o diagnóstico e a escolha do corretor com corridas de outros estudos e laboratórios, de preferência com mais de um lote e com mais QCs;
\item avaliar a escolha do corretor em dados com intensidade conhecida que não venham do gerador, como amostras com padrões adicionados, e procurar uma medida, calculável na matriz do usuário, que acompanhe a recuperação da intensidade verdadeira além da variação dos QCs;
\item repetir a avaliação com pesquisadores de metabolômica de outros perfis e com observação do uso;
\item validar as duas funcionalidades presentes no aplicativo em versão beta, a análise de multiverso e a predição e o enriquecimento de vias metabólicas, que ficaram fora do escopo deste trabalho.
\end{itemize}